\chapter*{缩写与读码约定}
\addcontentsline{toc}{chapter}{缩写与读码约定}

\begin{description}
  \item[LCQI] Local CPU Quantized Inference，本书第三册的本地 CPU 量化推理贯穿项目。
  \item[CLI] Command Line Interface，命令行入口。源码里通常只负责参数解析和结果打印，不承载核心推理逻辑。
  \item[MLP] Multi-Layer Perceptron，多层感知机。tiny MLP 用来建立可手算、可测试的最小推理闭环。
  \item[KV Cache] Key/Value Cache，decoder-only 模型跨 token 复用 attention key/value 的状态缓存。
  \item[BPE] Byte Pair Encoding，GPT-2 使用的 byte-level 子词编码方法。
  \item[GGUF] GPT-Generated Unified Format，llama.cpp 常用的模型文件格式。
  \item[TTFT] Time To First Token，服务请求从进入系统到输出第一个 token 的时间。
  \item[TPOT] Time Per Output Token，生成阶段每个输出 token 的平均或分位耗时。
  \item[SLO] Service Level Objective，服务等级目标。本书用它描述 TTFT、TPOT、拒绝率和取消回收等运行时约束。
  \item[PMU] Performance Monitoring Unit，硬件性能计数器单元。没有 PMU 权限时，性能报告不能伪装成硬件级结论。
  \item[AVX2] Advanced Vector Extensions 2，x86-64 SIMD 指令集扩展。源码存在不等于当前机器可用，报告要看运行时可用性。
\end{description}

读码时先看头文件中的数据结构和入口，再看实现文件中的检查、循环和错误分支，最后回到测试文件确认哪些行为已经被锁成回归合同。遇到 benchmark 或 smoke 脚本时，要区分“同模型同精度对比”和“成熟引擎参考口径”，不能把不同模型、不同量化格式的速度直接写成质量排名。
